% **************************************************************************************************************
% Université Gustave Eiffel QIRT-2022 template
%
% Author: Thibaud Toullier thibaud.toullier@univ-eiffel.fr
% Updates for QIRT Asia 2023: Andrei Sleptchenko andrei.sleptchenko@ku.ac.ae
% Updates for QIRT Asia 2027: QIRT Asia 2027 Technical Committee, program@qirtasia2027.com
%%
% License:
% This program is free software; you can redistribute it and/or modify
% it under the terms of the GNU General Public License as published by
% the Free Software Foundation; either version 2 of the License, or
% (at your option) any later version.
%
% This program is distributed in the hope that it will be useful,
% but WITHOUT ANY WARRANTY; without even the implied warranty of
% MERCHANTABILITY or FITNESS FOR A PARTICULAR PURPOSE.  See the
% GNU General Public License for more details.
%
% You should have received a copy of the GNU General Public License
% along with this program; see the file COPYING.  If not, write to
% the Free Software Foundation, Inc., 59 Temple Place - Suite 330,
% Boston, MA 02111-1307, USA.
%
% **************************************************************************************************************
% Note:
%    * To use the model : \documentclass{qirtconference}
%    * You can optionnaly set the `license` option to the class for loading CC license notice in the title page :
%      \documentclass[license]{qirtconference}
%    * You can optionnaly set the `license-nd` option to the class for loading CC license 
%      notice in the title page with No Derivative Works notice:
%      \documentclass[license-nd]{qirtconference}
%    * You can then customize the first page. Please note that you can left the arguments empty.
%      In that case, the line will not be displayed.
% **************************************************************************************************************
% To Do:
%    * Clean-up code, make sure it works
% **************************************************************************************************************


% Document class, specify license or not
\documentclass[license]{qirtconference}

\usepackage{lipsum}
% Title of your abstract / paper
\title{The title of your paper right here}

% List of authors
\author{A. Author}     % By default, one star will be added
\author[*]{B. Author}  % Equivalent to \author{B. Author}
\author[**]{C. Author} % Optionnaly: add stars to make corresponding affiliations

% Add the corresponding affiliations
\affiliation{Institute, Laboratory, Team, City, Country}
\affiliation[**]{Institute, Laboratory, Team, City, Country}

\begin{document}

% Create the title of the paper
\maketitle 

% write your abstract in the `abstract` environment
\begin{abstract}
This document is a template and already defines the principles components of your paper (title, text, heads, etc.).

This should not be more than \textbf{100 words}, no equation, no figure.
\end{abstract}

\section{Introduction}
This template is based on the previous MS Word and \LaTeX\  (by \href{mailto:thibaud.toullier@univ-eiffel.fr}{Thibaud Toullier}) templates and provides authors most of the formatting specifications needed for preparing, Extended Abstract and Conference Proceeding Paper.

\section{Writing a paper}

\subsection{Selecting a template}

First, you can choose the template that fits your requirements : MS Word or \LaTeX\ (this one). The \LaTeX\ template is also available on Overleaf, so you can use it as a template directly on your web browser for ease.
The paper size is A4 in any case.

\subsection{Maintaining the specifications}

The template is made to standardize the submitted papers, please do not alter the format and the style of text. Margins, line spaces and text fonts should not be altered. Please \textbf{do not add page numbering}.
However if you notice something that need to be improved, feel free to submit your recommendations and/or issues at \href{mailto:program@qirtasia2027.com}{program@qirtasia2027.com}

\subsection{Choosing a license}
You can choose between two licenses by default. For more information about the licenses, you can access the \href{https://creativecommons.org}{Creative Commons Website}. Available licenses:
\begin{itemize*}
\item CC BY default or with option « \verb~license~ » (\verb~\documentclass[license]{qirtconference}~)
\item CC BY ND with option « \verb~license_nd~ » (\verb~\documentclass[license_nd]{qirtconference}~)
\end{itemize*}

\section{Paper properties}

This is what the normal text looks like. This format is common for Extended Abstract as well as Conference Proceeding Paper.

\subsection{Extended Abstract}
Extended Abstract should be submitted for possible presentation in QIRT Asia 2027. It should be composed of Title, Author, Affiliation and Abstract (100 words maximum) with body text and \textbf{2 pages maximum}. The deadline is \textbf{15 October 2026}.

\subsection{Conference Proceeding Paper}
Conference Proceeding Paper should be submitted for publication in QIRT Asia 2027 if you are accepted for presentation. It should be composed of Title, Author, Affiliation and Abstract (100 words maximum) with body text and more than 6 pages but not exceed 10 pages. The deadline is \textbf{10 January 2027}.

\section{Heading 1}

This is an example of how the paper is rendered. The font use is the default \LaTeX\ sans-serif \verb~Computer Modern Sans~ font.

This is a new paragraph. As you can see, all paragraphs are indented, even the first one in the section.

\subsection{Equations}
Equations are given as follow and referred to in the text as Eq. \ref{eq:equation1} (\verb~Eq. \ref{eq:equation1}~) or Eqs. \ref{eq:equation1} and \ref{eq:equation2} ((\verb~Eqs. \ref{eq:equation1} and \ref{eq:equation2}~)). Note that it is preferred that all symbols are in italics both in equations and in the text. Please use the SI system for units (see \verb~siunitx~ package for instance).
\begin{equation}
  \nabla^{2}T - \frac{1}{\alpha}\frac{\partial T}{\partial t} = 0
  \label{eq:equation1}
\end{equation}

The same equation again:
\begin{equation}
  \nabla^{2}T - \frac{1}{\alpha}\frac{\partial T}{\partial t} = 0
  \label{eq:equation2}
\end{equation}

\subsubsection{Figures and tables}
All figures (i.e. graphs, data plots, illustrations, photographs, …) must be numbered in the order that they are introduced and referred to as figure \ref{fig:figure1} (\verb~figure \ref{fig:figure1}~). Authors may choose to insert the figures in the text as they occur (full or part width of the page) or to put them all at the end of the paper after references. In both cases, each figure should have title under it in the style shown below. To increase quality of figures the use of vector graphics with colors is recommended, but remember to adopt colors that remain clear if printed in black and white.

\begin{figure}[ht]
  \centering
  \includegraphics[width=0.3\linewidth]{example-image}
  \caption{\label{fig:figure1} Caption of a figure}
\end{figure}


Tables may also be in the text (full or part width of the page) or after references (and before the figures if placed at the end of the paper). The title of the table should be put on top of the table. Tables have to be referred to in the text as table \ref{table:table1} (\verb~table \ref{table:table1}~).

\begin{table}[h!]
\centering
\caption{\label{table:table1} Table to test captions and labels.}
\begin{tabular}{|c|c|c|c|}
 \hline
 Col1 & Col2 & Col2 & Col3 \\ [0.5ex] 
 \hline
 1 & 6 & 87837 & 787 \\ 
 \hline
 2 & 7 & 78 & 5415 \\
 \hline
 3 & 545 & 778 & 7507 \\
 \hline
 4 & 545 & 18744 & 7560 \\
 \hline
 5 & 88 & 788 & 6344 \\ [1ex] 
 \hline
\end{tabular}

\label{table:1}
\end{table}

\subsection{Lists}
List are accessible through the \verb~itemize~ environment. However, the standard spacing between the items may be too large. To reduce the spacing, we provide a \verb~itemize*~ environment:

\verb~itemize~ environment:

\begin{itemize}
\item First item in \verb~itemize~ environment
\item Second item in \verb~itemize~ environment
\end{itemize}

\verb~itemize*~ environment:

\begin{itemize*}
\item First item in \verb~itemize*~ environment
\item Second item in \verb~itemize*~ environment
\end{itemize*}

\subsection{References}
References are cited by order of appearance, like this \cite{MARTIN19771, Kaczmarek2007, inproceedings, RENKIELSKA2006867} with the \verb~unsrt~ bibliography style.

% Add your bibliography style here
\bibliography{refs}

\end{document}